% Chapter 4, Figure 6(b): percent error in y_b, B4 engine (scan: figures/ch4/fig06b.png).
% Curves: fig06.py (the book's interval method, eqs. (83)-(87) with the mass by eq. (74), against the
% Fehskens-Malewicki eqs. (20), (21) and the Caporaso-Bengen eqs. (27), (28)); read from fig06b.csv.
\documentclass[11pt]{standalone}
\usepackage{tamrfig}
\begin{document}
\begin{tikzpicture}
  \begin{axis}[tamr,
      width=4.0in, height=2.5in,
      xmin=0.02, xmax=0.10, xtick={0.02,0.04,...,0.10}, minor xtick={0.03,0.05,...,0.09},
      xticklabel style={/pgf/number format/.cd, fixed, fixed zerofill, precision=2},
      ymin=-15, ymax=15, ytick distance=5,
      xlabel={$m_o$ (kg)}, ylabel={Percent error in $y_b$},
      legend pos=north east, legend style={nodes={inner sep=1pt}}]
    \draw[muted, line width=0.4pt] (axis cs:0.02,0) -- (axis cs:0.10,0);
    \addplot[series1] table[x=m0, y=fm_kmax, col sep=comma] {figures/v2/ch4/fig06b.csv};
    \addlegendentry{Fehskens-Malewicki}
    \addplot[series2] table[x=m0, y=cb_kmax, col sep=comma] {figures/v2/ch4/fig06b.csv};
    \addlegendentry{Caporaso-Bengen}
    \addplot[series1] table[x=m0, y=fm_kmin, col sep=comma] {figures/v2/ch4/fig06b.csv};
    \addplot[series2] table[x=m0, y=cb_kmin, col sep=comma] {figures/v2/ch4/fig06b.csv};
    % k labels with a leader to each method's curve for that k (points on the fig06b.csv curves)
    \node (kmax) at (axis cs:0.033,-3.9) {\footnotesize $k_{\max}$};
    \draw[leader] (kmax.north east) -- (axis cs:0.038,-1.809);
    \draw[leader] (kmax.west) -- (axis cs:0.026,-4.364);
    \node (kmin) at (axis cs:0.027,-12.0) {\footnotesize $k_{\min}$};
    \draw[leader] (kmin.north) -- (axis cs:0.026,-5.337);
    \draw[leader] (kmin.north east) -- (axis cs:0.032,-8.943);
  \end{axis}
\end{tikzpicture}
\end{document}
